% 参考文献。注意：main.tex 中 \referencescn 已通过 thebibliography 自动生成
% “参考文献”一级标题（不带编号），并已 \addcontentsline 进入目录，
% 所以本文件**不要再写一级标题**，否则会出现两个“参考文献”标题。

\begin{thebibliography}{99}
\small
\setlength{\itemsep}{0pt}\setlength{\parsep}{0pt}\setlength{\parskip}{0pt}%
\bibitem{ref1} 姜启源, 谢金星, 叶俊. 数学模型[M]. 5 版. 北京: 高等教育出版社, 2018.
\bibitem{ref2} 司守奎, 孙玺菁. 数学建模算法与应用[M]. 3 版. 北京: 国防工业出版社, 2021.
\bibitem{ref3} Birge J R, Louveaux F. Introduction to Stochastic Programming[M]. 2nd ed. New York: Springer, 2011.
\bibitem{ref4} Hyndman R J, Athanasopoulos G. Forecasting: Principles and Practice[M]. 3rd ed. Melbourne: OTexts, 2021.
\bibitem{ref7} Hoerl A E, Kennard R W. Ridge regression: biased estimation for nonorthogonal problems[J]. Technometrics, 1970, 12(1): 55-67.
\end{thebibliography}
